% !TEX root = ../main_zh.tex
\chapter{菲利普斯曲线、预期与政策}
\label{ch:phillips}

总供给曲线告诉我们：只有当价格意外地高于人们的预期时，产出才会越过潜在水平。把这个意外翻译成通胀与失业的语言，它就成了菲利普斯曲线（Phillips curve）——通胀与失业之间那条短期的取舍菜单，而政策似乎正可以加以利用。但它究竟能不能被利用，完全取决于预期；而预期这个问题——公众如何形成预期、政策又能否塑造预期——正是现代货币与财政政策辩论的主线：政策该主动还是被动，该依规则还是凭相机抉择。本书的最后一章，便顺着这条线索从菲利普斯曲线一路走到泰勒规则。

\section{菲利普斯曲线}

从经验上看，在一个完整的经济周期里，通胀与失业总是朝相反的方向变动。现代的、引入了预期的菲利普斯曲线把这一关系写成
\[
\pi - \pi^{e} = -\beta\,(u - u^{n}) + v, \qquad \beta>0,
\]
其中 $u-u^{n}$ 是周期性失业（失业率对其自然率的偏离），$v$ 是供给冲击，$\pi^{e}$ 是预期通胀。当失业率低于自然率时，实际通胀就高于预期。

\defn{菲利普斯曲线}{
\emph{引入预期的菲利普斯曲线}（expectations-augmented Phillips curve）把意外通胀同周期性失业联系起来：$\pi-\pi^{e} = -\beta(u-u^{n}) + v$。它所刻画的取舍是通胀与\emph{周期性}失业之间的取舍，并且是相对于预期通胀而言的。}

\subsection{推导}

菲利普斯曲线其实就是改头换面的供给曲线。从短期总供给出发，$Y=\overline{Y}+\alpha(P-P^{e})$，解出价格水平，
\[
P = P^{e} + \frac{1}{\alpha}\,(Y-\overline{Y}).
\]
由于这条关系是从数据中读出来的，我们把供给冲击 $v$ 作为残差项附在后面，再对上一期的价格水平 $P_{-1}$ 取差分，
\[
P - P_{-1} = (P^{e}-P_{-1}) + \frac{1}{\alpha}\,(Y-\overline{Y}) + v,
\]
用通胀率写出来，便是 $\pi - \pi^{e} = \dfrac{1}{k\alpha}\,(Y-\overline{Y}) + v$，其中 $k$ 是一个比例常数。最后引用\emph{奥肯定律}（Okun's law）——这条经验规律说产出与失业总是同向变动：粗略地说，失业率每多上升一个百分点，大约对应损失两个百分点的 GDP。奥肯定律让我们可以用周期性失业替换产出缺口，
\[
\frac{1}{k\alpha}\,(Y-\overline{Y}) = -\beta\,(u-u^{n}),
\]
代回去就得到了菲利普斯曲线。这一关系是相关而非因果，但它把一条政策反复试图加以利用的联系给量化了出来。

\section{预期}

一切都系于 $\pi^{e}$。在信息不完美或价格黏性的世界里，人们是依照自己的预期行事的，于是整条菲利普斯曲线的位置都会随这一预期一起移动。关于预期如何形成，有两套理论，它们给出的政策结论截然不同。

\subsection{适应性预期}

\emph{适应性}（adaptive）预期是把过去外推得到的。最简单的一种设定取 $\pi^{e}=\pi_{-1}$，这就把菲利普斯曲线变成了通胀的一条运动方程，
\[
\pi = \pi_{-1} - \beta\,(u-u^{n}) + v.
\]
通胀因此具有惯性：在没有冲击的情况下，过去的通胀会传导到当前乃至未来的通胀。这本质上就是对历史数据做回归所得到的结果，而它的软肋也正在于此——它预设过去那些关系会一直延续下去。

\rmkb[卢卡斯批判]{用历史数据估计出来的关系去预测政策效果是不可靠的，恰恰因为政策本身会改变预期。用卢卡斯（Lucas）的话说：``Forecasting the effects of policy changes has often been done using models estimated with historical data. Such predictions would not be valid if the policy change alters expectations in a way that changes the fundamental relationships between variables.'' 当政策体制发生变化时，历史相关关系会出现\emph{结构断裂}（structural break），适应性预期的预测便随之失效。}

\subsection{理性预期}

\emph{理性}（rational）预期动用的是模型和一切可得的信息，而不只是过去——预期实际上是在预测自己以及所有其他人的决策。因此分析就不能只靠回归，它需要一个一般均衡模型；这个模型虽然未必能厘清单个的因果，却把每一条机制都纳入其中，给出一套自洽的冲击—反应（shock--response）。在理性预期下，原则上可以实现\emph{无痛去通胀}（painless disinflation）：如果中央银行能用一个可信的信号直接改变预期，供给曲线就会移动，通胀随之下降而不必经历衰退——牺牲率（sacrifice ratio）未必成为约束。对美国而言，理性预期约等于``相信美联储''。而当央行释放的信号过于晦涩、不足以支撑起量化的预期——政策的逻辑难以读懂——理性预期就退化成了小道消息，市场也随之频频震荡。

\section{自然率假说与延滞性}

这一切的背后是\emph{自然率假说}（natural-rate hypothesis）：存在一个潜在产出水平（以及一个自然失业率），经济在那里处于均衡，长期中必将到达、并且终究总会回归，而短期不过是对它的偏离。这是宏观经济学的组织性假设之一。

与之竞争的观点是\emph{延滞性}（hysteresis）：暂时性冲击可以推动长期水平本身，于是潜在产出虽然存在，却并非一个固定不变的数字——它取决于经济自身的历史。疫情就是一个现成的例子。一段长时间的失业会侵蚀工人的技能，等到需求复苏时他们已找不到与之匹配的工作，议价能力随之丧失；旷日持久的衰退会让消费习惯转向预防性（precautionary）储蓄；供应链外移，资本损耗，投资趋于保守。每一条渠道都让一个本应转瞬即逝的冲击留下了永久的印记。

\section{政策的实施}

\subsection{主动与被动}

\emph{主动}（active）的立场认为，在经济萧条时政府有责任采取行动——减轻民众的痛苦、刺激经济。但有两个问题让主动政策变得复杂。

第一，政策的作用存在\emph{时滞}（lags）。\emph{内部时滞}（inside lag）是从认识到需要行动、到真正出台对策所花的时间——要确认经济确实在收缩，得靠好几项指标共同印证，而执行（尤其是财政政策的执行）又需要协调。\emph{外部时滞}（outside lag）是政策影响到经济所花的时间，而这一影响还可能矫枉过正。第二，经济本身含有\emph{自动稳定器}（automatic stabilizers）——这些制度无需刻意行动、也几乎没有时滞，便能熨平周期。累进所得税在繁荣期会随额外收入抽取越来越高的份额，从而抑制乘数效应；失业保险在衰退期支撑收入、进而支撑消费；在浮动汇率下，强劲的出口繁荣会推升本币、从而给需求降温。

要把主动政策用好，政府必须尽早判断出经济周期所处的阶段——观察那些先于周期变动的领先指标（综合指数、投资、PMI、房地产资金到位量），并在环境变化到历史关系出现断裂时，改用结构模型而非简单外推。它还必须\emph{领先于}公众的预期：如果一个理性的公众早已料到某项政策，这项政策的效果就会被削弱。假设过去的数据显示货币增长会降低失业率，公众于是预期到了这一点；那么当中央银行扩大货币供应时，预期通胀会同步上升，失业率并不会下降——被预料到的政策是中性的，此时需要一项不同的、或者出人意料的政策。

\subsection{规则与相机抉择}

更深一层的选择，是政策究竟应当遵循一条公开的\emph{规则}（rule），还是凭\emph{相机抉择}（discretion）来制定。

\defn{规则与相机抉择}{
\emph{依规则}行事的政策，是依照一条明确、公开的公式来设定和调整的，因而可以被预期。\emph{凭相机抉择}行事的政策，则视具体情形逐案应对，无法被预期。}

最典型的规则是关于联邦基金利率的\emph{泰勒规则}（Taylor rule），它根据通胀和产出缺口来设定目标利率，
\[
r_{\mathrm{ff}} = 2 + 0.5\,(\pi - 2) - 0.5\cdot \mathrm{GDP\ gap}, \qquad
\mathrm{GDP\ gap} = 100\cdot\frac{\overline{Y}-Y}{\overline{Y}},
\]
其中 $2\%$ 被取作自然通胀率。注意这里的符号约定：缺口被定义为产出\emph{低于}潜在水平的\emph{亏空}，因此减去它就意味着——当产出\emph{高于}潜在水平时（繁荣，$Y>\overline{Y}$，缺口为负），规则会\emph{抬高}利率；当产出低于潜在水平时则\emph{压低}利率——这正是稳定经济的反应，等价于加上通常意义上的产出缺口 $(Y-\overline{Y})/\overline{Y}$。这是一条透明而可信的规则：利率由公式确定，具有独立性和可预测性，而不必去追问某次通胀的来源是什么。相比之下，中国追求让 $M_2$ 增速``匹配''名义 GDP 增速——这个目标同货币数量论赋予货币的逆周期角色配合得相当别扭，而它的``匹配''程度、以及它所依赖的对名义 GDP 增速的预测，本身又都是相机抉择的。

\state{支持规则的理由}{
有若干论据支持规则胜过相机抉择。决策者可能并不具备所需的专业知识；他们的判断可能建立在模糊或多变的原则之上，而这些原则又可能与公众的并不一致；他们还可能存在时间不一致，从而损害\emph{可信度}（credibility）——而预期、进而政策自身的有效性，恰恰仰赖于这一可信度。在相机抉择的决策者其能力与一致性都存疑的地方，一条设计良好的规则可以做得更好。}
